\chapter{Calling C}\label{app-nest}

A program sometimes needs a function that exists only in another
language: a routine of the C library, of the operating system, or of a
library written in C. ga68 can call such functions directly, and it can
also use variables of a C program. This appendix describes how. The
mechanism is a GNU extension; a68j and Algol 68 Genie cannot call C
(a68j has a similar mechanism for Java methods, described in its own
documentation).

\section*{Formal holes}

The construct that names a C function or variable is a \term{formal
hole}:
\begin{lstlisting}
nest C "name"
\end{lstlisting}
It is written where a value is expected and stands for that value. Its
type is not written in the hole; it is the type the context requires, so
a formal hole may only be written in a strong context
(Section~\ref{sec-contexts}), such as the value of a constant
declaration, an argument or the right-hand side of an assignment:

\example{appF/libc}

\noindent Each declaration says which type the C function has in Algol 68
terms. The compiler translates the type into the corresponding C type
and generates a call of the C function with that name. The name
\code|C| after \code|nest| is the language; the word \code|C| is not
reserved and can still name a type, an operator or a module elsewhere in
the program. The C library and the mathematical library are linked with
every program by the ga68 command, so their functions need nothing
else.

The notation goes back to the proposal for modules and separate
compilation of 1978~\cite{modules} (the same proposal as the modules of
Appendix~\ref{app-modules}). There, a formal hole marks a place in one
part of a program that another, separately compiled part fills in, so
that a program can be developed top-down; the proposal already allowed
the part to be written in another language, with an example in Fortran.
ga68 provides only the second use.

\section*{Types}

The types that can cross the border are the simple ones and those
built from them:

\begin{center}
\small
\begin{tabular}{ll}
\toprule
Algol 68 & C \\
\midrule
\code|void| (as a result) & \texttt{void} \\
\code|bool| & \texttt{bool} \\
\code|char| & a 32-bit integer (the character code) \\
\code|int|, \code|short int|, \code|short short int| & \texttt{int}, \texttt{short}, \texttt{char} \\
\code|long int|, \code|long long int| & \texttt{long}, \texttt{long long} \\
\code|bits|, \code|short bits|, \code|short short bits| & \texttt{unsigned int}, \texttt{unsigned short}, \texttt{unsigned char} \\
\code|long bits|, \code|long long bits| & \texttt{unsigned long}, \texttt{unsigned long long} \\
\code|real|, \code|long real| & \texttt{float}, \texttt{double} \\
\code|ref| \emph{T} & a pointer to \emph{T} \\
\code|proc| with parameters and result of these types & a function, or a pointer to one \\
a record with fields of these types & the corresponding \texttt{struct} \\
\code|string|, only as a parameter & three parameters, see below \\
\bottomrule
\end{tabular}
\end{center}

\noindent Other arrays and unions cannot be passed.

\begin{pitfall}
Nothing checks that the declared type matches the C function. The
compiler trusts the declaration, and the C library has no type
information to compare it with. A wrong type gives wrong results, not an
error message. The most likely mistake is the floating-point type: since
\code|real| is a C \texttt{float}, a function of the C library that takes
and returns \texttt{double}, such as \texttt{cbrt} or \texttt{sin}, must be
declared with \code|long real|, as in the example above; declared with
\code|real|, it receives and returns meaningless bit patterns. The
\texttt{float} variants have an \texttt{f} at the end of their name
(\texttt{cbrtf}).
\end{pitfall}

\section*{Variables}

A C variable can be used in two ways. Without further marking, the hole
stands for the value of the variable at the moment the hole is
evaluated, here when the declaration is executed; later changes by the C
program are not seen:
\begin{lstlisting}
int limit = nest C "limit";
\end{lstlisting}
With an ampersand before the name, it stands for the variable itself,
and the Algol 68 declaration has a reference type: assignments to it
change the C variable, and the C program sees them.
\begin{lstlisting}
ref int calls = nest C "&calls";
\end{lstlisting}
Without the ampersand, a hole of a \code|ref| type stands for the value of
a C pointer variable, so the Algol 68 name refers to the C object the
pointer points to.

\section*{Strings and functions as parameters}

A C function that receives an Algol 68 string gets three parameters for
it: a pointer to the characters, the number of characters, and the
distance in bytes from one character to the next. The characters are
32-bit codes, not bytes, and they need not be next to each other, because
the string may be a column of a two-dimensional array of characters,
such as \code|m[, 2]|. A C function that counts the
vowels of a string, a C function that calls a function it receives, and a
counter variable:

\cprogram{appF/cfun}

\noindent and the Algol 68 program that uses them:

\example{appF/cfun}

\noindent The C file is compiled with a C compiler and linked with the
program:
\begin{lstlisting}[style=shell]
$ cc -c cfun.c
$ ga68 -I. cfun.a68 cfun.o genietransput.o -o cfun
\end{lstlisting}
An Algol 68 function that is passed to C, such as \code|square| above,
arrives as an ordinary C function pointer. ga68 makes it point to a small
piece of code that it places in the stack frame of the block in which
the function is declared. A C library must therefore not keep the
pointer and call it after that block has ended, even if the function
uses no names of the block. The same piece of code is the reason why the
linker warns that such a program ``requires executable stack''; the
warning is expected, but a system that forbids executable stacks will
not run the program.
